\uuid{Vvxb}
\exo7id{3638}
\titre{exo7 3638}
\auteur{quercia}
\organisation{exo7}
\datecreate{2010-03-10}
\isIndication{false}
\isCorrection{false}
\chapitre{Endomorphisme particulier}
\sousChapitre{Autre}
\cpge{mpsi-mp2i.espaces-vectoriels.a, pcsi.espaces-vectoriels.a, ptsi.geometrie-espace.c, mpsi-mp2i.dimension-finie.a, pcsi.dimension-finie.a, ptsi.dimension-finie.a, mpsi-mp2i.matrices-applications-lineaires.a, pcsi.matrices-applications-lineaires.a, ptsi.matrices-applications-lineaires.a, bcpst1.matrices, bcpst1.espace-vectoriel-kn.b, bcpst2.espaces-vectoriels.a, bcpst2.espaces-vectoriels.b, bcpst2.applications-lineaires-matrices.c}
\module{Algèbre}
\niveau{L2}
\difficulte{}

\contenu{
\texte{
Soient $f_1,\dots,f_n : \R \to \R$ $n$ fonctions linéairement indépendantes.
On pose $E = \text{vect}(f_1,\dots,f_n)$ et pour $x \in \R$,
$ {\delta_x} : E \to \R, f \mapsto {f(x).}$
}
\begin{enumerate}
    \item \question{Montrer que la famille $(\delta_x)_{x\in\R}$ engendre $E^*$.}
    \item \question{En déduire qu'il existe $x_1,\dots,x_n \in \R$ tels que $\det M \ne 0$ où
    $M$ est la matrice de terme général $f_i(x_j)$.}
\end{enumerate}
}
